\scanpage{350}
% top centre
\mspage{266}
% lines 1–7
changement de base soient des foncteurs images directes pour des morphismes de paratopos (i.e. des foncteurs accessibles, commutant aux \inserted{petites} $\varprojlim$). Supposons de plus $J$ équivalente à une petite catégorie. Alors $\underline{\unsure{\Gamma|\mathscr T}}(N/J)$ est une catégorie paratopos.

% lines 8–16
En effet, l'hypothèse implique que pour \deleted{\unsure{tout|toute} \unsure{$\lambda\longrightarrow$|de}} flèche $\lambda:j'\unsure{\mathrel{-}|\longrightarrow}j$ dans $J$, le foncteur $\lambda_*:N_{j'}\mathrel{-}N_j$ admet un adjoint à gauche $\lambda^*:N_j\mathrel{-}N_{j'}$, d'où s'ensuit que $N\unsure{\longleftarrow|\longrightarrow}J$ est aussi \unsure{cofibrant|cofibré}, les \deleted{\unsure{morphis}} foncteurs de \unsure{cochgt.} de base \unsure{commutant|commutent} aux petites $\unsure{\varinjlim|\underline{\lim}}$. Donc on peut appliquer 4.\deleted{\unsure{8}}\unsure{10|8}.9, qui donne le résultat voulu.

% lines 17–21
\uline{Question} 4.10.11. Soit $N\unsure{\longleftarrow|\mathrel{-}}J$ un foncteur \uline{fibrant}, tel que les $N_j$ soient des catégories accessibles, et les foncteurs $\lambda^*$ de \unsure{chgt.} de base accessibles, $J$ étant \unsure{équivalente|équivalent} à une petite catégorie. \unsure{Je|J-}
% archive sheet number: 349
